\documentclass[11pt,a4paper]{article}

\usepackage[margin=2.2cm]{geometry}
\usepackage{graphicx}
\usepackage{booktabs}
\usepackage{tabularx}
\usepackage{float}
\usepackage{setspace}
\usepackage{titlesec}
\usepackage[table]{xcolor}
\usepackage{caption}
\usepackage{enumitem}
\usepackage[T1]{fontenc}
\usepackage{lmodern}
\usepackage{microtype}
\usepackage[hidelinks,colorlinks=true,linkcolor=linknavy,citecolor=linknavy,
            urlcolor=linknavy]{hyperref}

\definecolor{linknavy}{HTML}{16447A}
\definecolor{tablehead}{HTML}{E4E9EF}

\setstretch{1.05}
\setlist{nosep,leftmargin=1.4em}
\titleformat{\section}{\large\bfseries}{\thesection.}{0.6em}{}
\titleformat{\subsection}{\normalsize\bfseries}{\thesubsection.}{0.6em}{}
\titlespacing*{\section}{0pt}{12pt}{5pt}
\titlespacing*{\subsection}{0pt}{8pt}{3pt}
\captionsetup{font=small,labelfont=bf,skip=4pt}
\setlength{\textfloatsep}{10pt}\setlength{\intextsep}{10pt}
\renewcommand{\arraystretch}{1.25}
\setlength{\parskip}{3pt}

\newcommand{\shot}[3]{%
  \begin{figure}[H]\centering
  \includegraphics[width=#2\linewidth]{figures/#1.png}
  \caption{#3}\label{fig:#1}\end{figure}}

\input{measurements.tex}

\newcommand{\authorA}{Munsif}
\newcommand{\authorB}{Ashfaq}
\newcommand{\authorC}{Lareef}
\newcommand{\submissiondate}{\today}
\newcommand{\repourl}{https://github.com/munsif-dev/hospital-vitals-kappa}

\begin{document}

\begin{titlepage}
    \centering
    \vspace*{0.5cm}
    \includegraphics[width=4cm]{Logo.jpeg}\par
    \vspace{1cm}
    {\LARGE Department of Electrical and Information Engineering}\\[0.3cm]
    {\Large Faculty of Engineering, University of Ruhuna}\par
    \vfill
    {\Large \textbf{Mini Project Report}}\\[0.3cm]
    {\large{EC8203 Applied Big Data Engineering}}\par
    \vfill
    {\huge\textbf{Hospital Patient Vital Signs Monitoring}}\\[0.5cm]
    {\large\itshape A Kappa architecture for watching a ward in near real time}\par
    \vfill
    \begin{tabular}{ll}
      \large EG/2021/4417 & \large ASHFAQ M.R.M \\[0.25cm]
      \large EG/2021/4644 & \large LAREEF J.M. \\[0.25cm]
      \large EG/2021/4684 & \large MUNSIF M.F.A \\
    \end{tabular}\\[0.8cm]
    {\large Computer Engineering}\par
    \vfill
    {\large \submissiondate}\\[0.4cm]
    {\small Source code: \url{\repourl}}\par
    \vspace*{0.5cm}
\end{titlepage}

\hypersetup{linkcolor=black}
\tableofcontents
\hypersetup{linkcolor=linknavy}
\clearpage

% =====================================================================
\section{The problem}
% =====================================================================

A hospital ward has 40 beds. Each bed has a monitor that measures heart rate,
oxygen level (SpO2), breathing rate, blood pressure and temperature. Once a day the
pathology lab sends a file with blood test results for every patient. The ward wants
to answer one question, taken from the assignment:

\noindent\textbf{Business question:} Which patients show concerning vital-sign trends
right now, and how do yesterday's lab results change the risk picture for those
patients going forward?

Nurses already use a standard score for this, the National Early Warning Score 2
(NEWS2)~\cite{rcp}. Each vital sign gets 0 to 3 points depending on how far it is from
normal, and the points are added up. A total of 5 or more means ``call a doctor
soon''; 7 or more means ``emergency''. Our system calculates NEWS2 for every reading,
adds points for abnormal blood tests, and shows the ward with the sickest patient at
the top.

\subsection{The two data sources}

Both sources are Python programs that simulate the real thing.

\begin{itemize}
  \item \textbf{Streaming source: bedside monitors.} One reading per bed every 15
  simulated minutes, about \mThroughput{} readings per real second for the whole ward.
  The numbers drift the way real vital signs do (each value wanders around the
  patient's own normal level, oxygen falling pushes the heart rate up, and so on).
  About 1.8\,\% of readings are deliberately broken, like a probe falling off, so we can
  check that the system catches them.
  \item \textbf{Daily source: the pathology lab.} One JSON file per simulated day at
  06:00 with six tests per patient (lactate, white cells, CRP, creatinine, haemoglobin,
  potassium). One day's file is missing, one has a corrupted reference range, and some
  arrive late, so that every branch of the ingest workflow is exercised.
\end{itemize}

\subsection{Simulated time}

One simulated day lasts 5 real minutes, so the clock runs 288 times faster than real
time. One real minute is 4.8 simulated hours. All containers read the same start time
from one small file, so they agree on the simulated time. The run starts on 1 April 2026.

\subsection{Three scripted patients}

Most patients are stable. Three follow a script, so that we know in advance what the
system should do with them:

\begin{itemize}
  \item \textbf{P014} develops sepsis from day 2, 08:00. Temperature, heart rate and
  breathing rate climb and blood pressure falls. The day-2 blood tests agree (lactate
  3.8, white cells 18.4). \emph{The system must raise an alert.}
  \item \textbf{P031} has COPD, a lung disease. Their normal oxygen level is 89 to
  91\,\%, which is low for most people but safe for them. \emph{This patient is used to
  test the replay (Section~\ref{sec:replay}).}
  \item \textbf{P007} has one heart-rate reading of 145 on day 1 at 14:00, then returns
  to normal. \emph{The system must not raise an alert for one odd reading.}
\end{itemize}

% =====================================================================
\section{What the system must do}
% =====================================================================

Table~\ref{tab:req} turns the use case into concrete requirements and shows where each
one is met.

\begin{table}[H]\centering\small
\begin{tabularx}{\linewidth}{|l|X|X|}
\hline
\rowcolor{tablehead}& \textbf{Requirement} & \textbf{How it is met} \\ \hline
R1 & See a worsening patient within seconds, not at the end of the day &
One stream scores every reading. Measured end-to-end time: \mLatencyPfifty{} s median. \\ \hline
R2 & Combine the daily lab file with the live vitals & Labs are published to a Kafka
topic and joined to each reading as it is scored. \\ \hline
R3 & Never lose a reading, and never score an impossible one & Impossible readings go
to a dead-letter topic with their original bytes; late ones go to a late topic. \\ \hline
R4 & One alert per problem, not one per reading & Alerts are grouped into episodes. \\ \hline
R5 & Be able to change the scoring rule and re-score the past & Kafka keeps 30
simulated days; a second copy of the stream re-reads them with the new rule. \\ \hline
R6 & A daily report and a live dashboard & Airflow renders a PDF each simulated day;
Grafana and a REST API show the ward live. \\ \hline
R7 & Notice when the system itself stops & Prometheus alert
\texttt{WardMonitoringSilent}, and an Airflow health check. \\ \hline
\end{tabularx}
\caption{Requirements taken from the use case.}\label{tab:req}
\end{table}

% =====================================================================
\section{Architecture decision: Kappa, not Lambda}\label{sec:decision}
% =====================================================================

\subsection{The two options}

\hspace*{\parindent}\textbf{Lambda} would have two paths. A speed layer would score readings as they
arrive, and a batch layer would recompute everything from stored history each night.
A serving layer would merge the two answers.

\textbf{Kappa} has one path. Every reading goes through one stream. History is kept in
Kafka itself. To recompute the past, you start a second copy of the same stream at the
beginning of the log.

\subsection{Why we chose Kappa}

\hspace*{\parindent}\textbf{1. The score must be the same everywhere.} In Lambda the NEWS2 rule would
exist twice, once in the stream code and once in the batch code. If they ever differ,
even by one boundary (is a heart rate of 131 in the 2-point band or the 3-point band?),
the ward screen and the daily report would disagree about the same patient. In a
hospital that is not a small reporting error; it is a patient who looked fine on one
screen and unwell on another. With Kappa the rule exists once, in
\texttt{ward/clinical/news2.py}. A test reads the source code of the whole project and
fails if any other file contains a scoring table.

\textbf{2. The daily file is small reference data, not a big batch job.} Each day's lab
file has 240 rows. It does not need a separate batch engine; it needs to be joined to
the stream. We publish it to a compacted Kafka topic (the newest result per patient and
test is always kept), and the stream keeps it in memory for the join.

\textbf{3. The main reason to recompute history is a rule change.} Clinical rules
change. NEWS2 itself changed in 2017 to add a second oxygen scale for COPD patients.
When that happens we need to re-score past readings with the new rule and compare. In
Kappa this is a replay of the log through the same code. In Lambda it is a batch job,
which means writing the new rule twice.

\textbf{4. Latency.} Every reading must be scored within seconds. Both architectures
can do that, so latency does not decide between them. But it rules out any design where
the ward waits for a nightly job.

\subsection{What we give up}

Kappa is not free. We accept these costs:

\begin{itemize}
  \item \textbf{Replay is slower than a bulk batch job} for very long histories,
  because it goes through the stream one micro-batch at a time. For 30 simulated days
  of one ward it is fast (Section~\ref{sec:replay}); for years of data from a whole
  hospital it would not be.
  \item \textbf{History lives in Kafka.} Kafka must keep 30 days of readings on disk.
  A data lake would be cheaper per gigabyte.
  \item \textbf{Compacted topics forget old values.} The labs topic keeps only the
  newest result per test. A replay therefore sees fewer historical lab results than the
  live stream did. We compare the NEWS2 part of the score, which does not use labs, and
  we say so wherever it matters.
\end{itemize}

\subsection{When we would choose Lambda instead}

We would switch if the daily source became a heavy analytic job (for example a
nightly model trained on months of data from every ward), if history had to be kept for
years, or if recomputing the past quickly mattered more than having one copy of the
logic.

% =====================================================================
\section{How the system is built}
% =====================================================================

Figure~\ref{fig:d1} shows who uses the system. Figure~\ref{fig:d2} shows every
service, the port it uses, and how data moves from top to bottom.

\begin{figure}[H]\centering
\includegraphics[width=0.8\linewidth]{figures/D1-context.pdf}
\caption{System context: two data sources in, four kinds of user out.}\label{fig:d1}
\end{figure}

\begin{figure}[H]\centering
\includegraphics[width=1.0\linewidth]{figures/D2-layered-architecture.pdf}
\caption{Layered architecture. Blue is streaming, orange is storage and serving, green
is orchestration, grey is observability. Every service runs in Docker Compose.}
\label{fig:d2}
\end{figure}

\subsection{Ingestion: Kafka}

Kafka is the system of record. Table~\ref{tab:topics} lists the six topics. The
cleanup policy matters: readings are deleted after 30 simulated days, while the
admissions and labs topics are \emph{compacted}, so the newest value for each key is
kept forever. Every message is Avro, checked against a schema in the Schema Registry,
and keyed by patient, so one patient's readings always arrive in order.

\begin{table}[H]\centering\small
\begin{tabularx}{\linewidth}{|l|l|l|X|}
\hline
\rowcolor{tablehead}\textbf{Topic} & \textbf{Partitions} & \textbf{Cleanup} & \textbf{What it holds} \\ \hline
vitals.readings.v1 & 6 & delete after 30 sim days & every bedside reading \\ \hline
labs.results.v1 & 3 & compact & newest result per patient and test \\ \hline
ward.admissions.v1 & 3 & compact & bed, age, condition, COPD flag \\ \hline
vitals.alerts.v1 & 3 & delete & new alert episodes \\ \hline
vitals.readings.dlq & 1 & delete & impossible readings, with original bytes \\ \hline
vitals.late & 1 & delete & readings that arrived too late for their window \\ \hline
\end{tabularx}
\caption{Kafka topics (measured with \texttt{make topics}).}\label{tab:topics}
\end{table}

The lab file reaches Kafka through the Airflow workflow \texttt{ward\_lab\_ingest}. It
checks the file's SHA-256 checksum, checks every row, moves a bad file to a quarantine
folder, and otherwise publishes each row unchanged.

\subsection{Processing: one Spark stream}

The stream is a Spark Structured Streaming job. Every 5 real seconds it takes the new
readings from Kafka and handles each one as shown in Figure~\ref{fig:d6}.

\begin{figure}[H]\centering
\includegraphics[width=0.45\linewidth]{figures/D6-one-reading.pdf}
\caption{What happens to one reading.}\label{fig:d6}
\end{figure}

The steps are:

\begin{enumerate}
  \item \textbf{Check.} A reading that cannot be decoded, or that is physically
  impossible (SpO2 of 0, heart rate of 300, systolic pressure below diastolic), goes to
  the dead-letter topic with its original bytes and its exact position in Kafka. We use
  physical limits, not statistics. A statistical filter would throw away the most
  important readings, because a patient who is getting worse is, by definition, far
  from average.
  \item \textbf{Duplicates and late readings.} A reading id seen before is skipped. A
  reading more than 60 simulated minutes older than the newest one for that patient
  goes to the late topic and is stored, but not scored.
  \item \textbf{Trend.} The last 4 simulated hours (16 readings) are kept per patient,
  and a slope is fitted to each vital sign. We use the Theil-Sen slope (the median of
  all pairwise slopes), because one odd reading cannot move it much.
  \item \textbf{Score.} NEWS2 is calculated by \texttt{news2.py}. Patients with the
  COPD flag use oxygen Scale 2 in version 2 of the rule.
  \item \textbf{Labs.} Lab results reported \emph{before} the reading add up to 5
  points (for example +1 for lactate above 2.0). This gives the composite risk.
  \item \textbf{Alerts.} Seven rules are checked, from ``NEWS2 7 or more'' to ``two
  vital signs trending the wrong way''. An alert is raised only when a new episode
  starts or its severity rises, and a single extreme reading must be confirmed by the
  next one.
\end{enumerate}

Figure~\ref{fig:d3} follows patient P014 through these steps.

\begin{figure}[H]\centering
\includegraphics[width=0.62\linewidth]{figures/D3-event-sequence.pdf}
\caption{Event sequence for patient P014 on simulated day 2.}\label{fig:d3}
\end{figure}

\subsection{Storage: Cassandra}

Scores land in Cassandra. We designed one table for each question the system has to
answer, so every read goes to exactly one partition and no query needs
\texttt{ALLOW FILTERING} (Figure~\ref{fig:d5}). Three choices matter most:

\begin{itemize}
  \item \texttt{ward\_risk\_snapshot} is ordered by risk score inside the partition,
  so the ward list comes back sickest-first without sorting.
  \item It has a 120-second time to live. If the stream stops, the ward list
  \emph{empties} within two minutes. An empty screen is an obvious failure; a screen
  frozen on old values would look like a calm ward.
  \item Every score row carries its \texttt{scorer\_version}, so version 1 and version 2
  results sit side by side during a replay and nothing is overwritten.
\end{itemize}

\begin{figure}[H]\centering
\includegraphics[width=0.72\linewidth]{figures/D5-cassandra-tables.pdf}
\caption{Cassandra tables and the query each one serves.}\label{fig:d5}
\end{figure}

\subsection{Serving and orchestration}

A FastAPI service reads Cassandra and returns the ward list, a patient's history, labs,
alerts, and the version comparison. Grafana shows three dashboards. Airflow runs five
workflows (Table~\ref{tab:dags}). Airflow only starts, checks and waits for things; it
never calculates a score. A test fails if a workflow imports a data-processing library.

\begin{table}[H]\centering\small
\begin{tabularx}{\linewidth}{|l|l|X|}
\hline
\rowcolor{tablehead}\textbf{Workflow} & \textbf{Runs} & \textbf{What it does} \\ \hline
ward\_lab\_ingest & every 5 min (1 sim day) & checksum, validate, quarantine or publish the lab file \\ \hline
ward\_daily\_report & every 5 min & wait for yesterday's summaries, render the PDF report \\ \hline
ward\_healthcheck & every 2 min & fail if readings stop, the ward list is short, or the stream stalls \\ \hline
ward\_retention & hourly & check the time-to-live settings and that the snapshot stays small \\ \hline
ward\_replay & on demand & wait for the replay, compare versions, stop for human approval \\ \hline
\end{tabularx}
\caption{Airflow workflows.}\label{tab:dags}
\end{table}

% =====================================================================
\section{Why each technology}
% =====================================================================

\begin{table}[H]\centering\small
\begin{tabularx}{\linewidth}{|l|X|X|}
\hline
\rowcolor{tablehead}\textbf{Tool} & \textbf{Why it fits this use case} & \textbf{Alternative we considered} \\ \hline
Apache Kafka (KRaft) & Keeps the history for replay; partitions by patient keep
each patient's readings in order; compaction suits the labs and admissions. &
RabbitMQ: deletes messages once read, so no replay. \\ \hline
Schema Registry, Avro & A reading that does not match the contract is rejected at the
door, not deep inside the stream. & JSON without a schema. \\ \hline
Spark Structured Streaming & Required by the brief. Micro-batches every 5 s are far
faster than the 15-minute reading interval. & Apache Flink: lower latency we do not
need. \\ \hline
Apache Cassandra & Writes are cheap; the ward list is one partition read; time to
live empties the screen if the stream stops. & PostgreSQL: fine at this size, but
sorting and expiry would need extra work. \\ \hline
FastAPI & Typed responses and an automatic documentation page. & Flask. \\ \hline
Apache Airflow & Required by the brief. Good at ``wait, check, branch, retry'' for
daily files. & cron: no retries, no history. \\ \hline
Prometheus, Alertmanager, Grafana & Pull-based metrics, alert rules as code, and
routing of patient-safety alerts separately from engineering alerts. & Logs only. \\ \hline
\end{tabularx}
\caption{Technology choices.}\label{tab:tech}
\end{table}

% =====================================================================
\section{Changing the scoring rule: the replay}\label{sec:replay}
% =====================================================================

This is the feature that justifies Kappa, so we tested it end to end. Version 1 of the
rule scores every patient's oxygen on Scale 1. Version 2 uses Scale 2 for patients with
the COPD flag. We wrote down the correct result before running it: \emph{patients
without COPD must not change at all}. For COPD patients, scores should fall when their
usual oxygen level of 88 to 92\,\% is no longer counted as abnormal. Our first version
of the check also said ``and never rise''. The first replay proved that wrong, and
rightly so: Scale 2 also gives points to a COPD patient on oxygen whose level is 93\,\%
or more, because too much oxygen is harmful for them. Every upward change we saw was
exactly that case, so the check now allows it and fails on anything else.
Figure~\ref{fig:d4} shows the steps.

\begin{figure}[H]\centering
\includegraphics[width=0.58\linewidth]{figures/D4-replay-branching.pdf}
\caption{Replay and cutover.}\label{fig:d4}
\end{figure}

\texttt{make replay} starts \texttt{ward-stream-v2}. It is the same image and the same
code as the live stream, with three settings changed: the rule version, the starting
position (the beginning of the log) and its own checkpoint. We measured:

\begin{table}[H]\centering\small
\begin{tabular}{|l|r|}
\hline
\rowcolor{tablehead}\textbf{Measurement} & \textbf{Result} \\ \hline
Readings in the log at the start of the replay & \mReplayEvents{} \\ \hline
Time for the v2 stream to catch up (real seconds) & \mReplaySeconds{} \\ \hline
Readings scored by both versions & \mReplayCompared{} \\ \hline
Readings whose NEWS2 changed & \mReplayChanged{} \\ \hline
\quad of which COPD patients & \mReplayCopdChanged{} \\ \hline
\quad of which other patients (must be 0) & \mReplayControlChanged{} \\ \hline
Changes that went up (COPD, on oxygen at 93\,\% or more) & \mReplayUpward{} \\ \hline
\quad of which not explained by that (must be 0) & \mReplayUnexplained{} \\ \hline
Readings that drop below NEWS2 5 under v2 & \mReplayBelowFive{} \\ \hline
P031: emergency (NEWS2 7 or more) alerts under v1 and v2 & \mReplayPthirtyoneAlerts{} \\ \hline
\end{tabular}
\caption{Replay results, from \texttt{make diff}.}\label{tab:replay}
\end{table}

P031's emergency alerts disappear under v2, but its total number of alerts went up
(27 under v1, 34 under v2). Under v1 its oxygen level kept it ``red'' all the time, which is
one long alert episode. Under v2 its score sits near the thresholds and moves across
them, and every crossing opens a new episode. A real system would add a margin before
an episode closes; we list this as a limitation.

The Airflow workflow \texttt{ward\_replay} runs the same comparison and ends at an
approval step. Airflow does not switch versions itself. The switch is one API call
(\texttt{make cutover VERSION=v2}); we tested it and the ward list changed from P031 at
risk 8 to risk 5, with P014 unchanged at 10. Going back is the same call with ``v1''.

% =====================================================================
\section{Observability}
% =====================================================================

A monitoring system that fails silently is dangerous: an empty alert list looks exactly
like a calm ward. So we measure the pipeline as carefully as the patients.

\begin{table}[H]\centering\small
\begin{tabularx}{\linewidth}{|p{3.2cm}|X|X|}
\hline
\rowcolor{tablehead}\textbf{What} & \textbf{How} & \textbf{Why} \\ \hline
Structured logs & JSON lines from every service, each with \texttt{service} and
\texttt{stage} (ingest, process, store, serve, orchestrate). & One query finds every
log line for one stage across all containers. \\ \hline
Readings in, scores out & Counters \texttt{readings\_produced\_total},
\texttt{risk\_scores\_written\_total}. & If scores stop while readings continue, the
stream is dead. \\ \hline
End-to-end time & Histogram from the monitor's timestamp to the Cassandra write. &
This is what a nurse waits for. \\ \hline
Micro-batch progress & A Spark query listener records batch time, rows and last
progress. & Tells a stalled query from an idle one. \\ \hline
Injected vs dead-lettered & Counters by fault type on both sides. & The injector is a
control: what goes in must come out. \\ \hline
Replay progress & Per-partition offsets of the replay stream. & Spark does not use
consumer groups, so consumer lag cannot measure it. \\ \hline
\end{tabularx}
\caption{What is measured, how, and why.}\label{tab:obs}
\end{table}

Seven Prometheus rules watch these signals. The most important is
\texttt{WardMonitoringSilent}: no scores written for one minute. It is written so that
it still fires when the stream container has disappeared completely, not only when it is
idle. Alertmanager sends it to two receivers, the ward and the on-call engineer, because
it is both a patient-safety and a technical problem.

We tested it by stopping the stream (\texttt{make chaos-stop-stream}). The alert moved
to pending after \mSilentPending{} s and fired after \mSilentFiring{} s, and it resolved
\mSilentResolved{} s after the stream was restarted (half of that is Spark starting).
No reading was lost: the stream restarted from its checkpoint and worked through the
backlog, and the day's summary carried on from what was stored (99 readings per patient
for that day, the same as the replay stream counted). A second rule, \texttt{StaleLabDataElevated}, fired on its own later in the run,
when the simulated lab withheld the day-5 file and the day-4 file was quarantined, so no
patient had lab results newer than two days. A test checks that every metric
named in an alert rule or a dashboard is really produced by the code; the first version
of the project had an alert on a metric that did not exist, so it could never fire.

We did not add distributed tracing. Each dead letter carries the exact Kafka partition
and offset of its reading, which is enough to find it, but a trace across all services
would be the next step.

% =====================================================================
\section{Results}
% =====================================================================

We ran the whole stack from nothing (\texttt{make clean}, \texttt{make up}) and measured
it. Table~\ref{tab:results} gives the numbers; the screenshots below and in
Appendix~\ref{app:shots} are from that run.

\begin{table}[H]\centering\small
\begin{tabularx}{\linewidth}{|X|r|}
\hline
\rowcolor{tablehead}\textbf{Measurement} & \textbf{Result} \\ \hline
Readings delivered to Kafka per second (expected about 12.8) & \mThroughput{} \\ \hline
Distinct patient keys on the vitals topic (expected 40) & \mKeys{} \\ \hline
Impossible readings injected (snapshot) & \mInjected{} \\ \hline
Of those, found in the dead-letter topic (Table~\ref{tab:dlq}) & \mDeadLettered{} \\ \hline
End-to-end time, median and 95th percentile (real seconds) & \mLatencyPfifty{} / \mLatencyPninetyfive{} \\ \hline
Micro-batch duration, typical (real seconds) & \mBatchSeconds{} \\ \hline
P014: first NEWS2 7+ alert after sepsis onset (simulated time) & \mPfourteenDelay{} \\ \hline
P014: NEWS2 7+ alerts raised in that episode (expected 1) & \mPfourteenHighAlerts{} \\ \hline
P007: alerts within 2 simulated hours of the spike (expected 0) & \mPsevenAlerts{} \\ \hline
API ward list response time, median & \mApiMs{} ms \\ \hline
Memory used by the whole stack & \mMemory{} \\ \hline
Automated tests passing & \mTests{} \\ \hline
\end{tabularx}
\caption{Measured results.}\label{tab:results}
\end{table}

Table~\ref{tab:dlq} checks the dead-letter path against the fault injector, reason by
reason. We read the injector's counter, waited 20 seconds, and counted the dead-letter
topic by distinct Kafka position. Every injected fault was there under its own reason,
and no other reason appeared, so no real reading was rejected. (In an earlier run this
check found 16 extra ``inverted blood pressure'' rejections: real P014 readings, made
impossible by a bug in the sepsis script. That bug was fixed before this run.)

\begin{table}[H]\centering\small
\begin{tabular}{|l|r|r|}
\hline
\rowcolor{tablehead}\textbf{Reason} & \textbf{Injected} & \textbf{In the dead-letter topic 20 s later} \\ \hline
SpO2 probe detached & 124 & 128 \\ \hline
Heart rate out of range & 86 & 87 \\ \hline
Empty reading & 75 & 78 \\ \hline
Blood pressure inverted & 51 & 51 \\ \hline
Temperature out of range & 50 & 51 \\ \hline
Timestamp in the future & 20 & 20 \\ \hline
Total & 406 & 415 \\ \hline
\end{tabular}
\caption{Injected faults against dead letters. The right column is slightly higher
because the monitors kept injecting during the 20 seconds. Nine duplicate readings were
also injected; they are skipped by id, not dead-lettered.}\label{tab:dlq}
\end{table}

\shot{grafana-ward-monitor}{1.0}{Grafana ward monitor during the run. The table lists
all 40 beds with the highest composite risk first. The chart shows the three scripted
patients: P014 rising on day 2, P031 steady, P007 with one spike.}

\shot{api-ward-monitor}{0.55}{The same ward list from the REST API
(\texttt{/api/v1/ward/WARD-A/monitor}), sickest first.}

\shot{cassandra-p014-alerts}{0.95}{Real \texttt{cqlsh} output: today's alerts. Each
condition appears once per episode, not once per reading.}

\shot{daily-report}{0.62}{First page of the daily PDF report produced by Airflow from
the stored daily summaries.}

\shot{grafana-replay}{1.0}{Replay dashboard after the v2 stream caught up. P031 (COPD)
scores lower under v2; P014 is identical under both.}

\shot{prometheus-alerts}{0.95}{Prometheus during the outage test:
\texttt{WardMonitoringSilent} firing after the stream was stopped.}

% =====================================================================
\section{Problems we found by checking against known answers}
% =====================================================================

Many faults in this project raised no error at all. They were found by comparing an
output with a value we knew in advance: 40 patient keys, the injector's own fault counts, zero
changes for non-COPD patients, one alert for P014. The most instructive ones:

\begin{itemize}
  \item No process could connect to Cassandra: the connection code passed an argument
  the driver does not accept. Every unit test replaced the database with a fake, so it
  only showed up on the first real run.
  \item The same alert was raised again every 15 simulated minutes, because its id was
  built from a time that moved with every reading. Alerts are now grouped in episodes.
  \item After the outage test, readings stamped two hours in the future were accepted,
  because the check compared them with the clock at processing time, which was far
  ahead while the stream caught up. It now compares with the time the reading was sent.
  \item P007's single spike was written into the patient's state, so the next two
  readings stayed high; later, a small clock jump put two readings in its time window.
  The monitor now sends it exactly once.
  \item Lab results were stored as ``LACTATE'' but looked up as ``lactate'', so labs
  never added any points. The live join also only read the labs topic once, at start-up.
  \item The dead-letter rate was almost double the injected rate, because oxygen values
  could drift above 100\,\% and were rejected as impossible.
  \item The API started ``healthy'' without a database connection, because an error
  on start-up was caught and ignored.
  \item After a restart in the middle of a day, the stream overwrote that day's
  summary with the few readings it had seen since, and joined the backlog with the wrong
  day's labs.
  \item The injector check found 16 real P014 readings rejected as impossible: the
  sepsis script lowered systolic pressure but not diastolic.
  \item The alert on the most important condition, a silent stream, named a metric that
  nothing produced, and most dashboard panels fell back to made-up constants.
  \item Most services printed plain-text logs although JSON was configured, because
  their loggers were created before the configuration was applied.
\end{itemize}

Each fix came with a test that failed before it. Where a test guards against a whole
class of mistakes (only one scoring rule, every alert names a real metric), we planted a
mistake to prove that the test catches it.

% =====================================================================
\section{Limitations and what we would do at production scale}
% =====================================================================

\begin{itemize}
  \item \textbf{One machine.} Kafka has one broker and Cassandra one node, with no
  replication. In production each would have at least three.
  \item \textbf{Scoring is done record by record in Python.} Spark reads the
  micro-batch, but the rows are collected to the driver and scored one at a time. This
  is fine for 13 readings per second; for a whole hospital the scoring would have to run
  inside Spark's workers.
  \item \textbf{Stream state is in memory.} The daily totals are reloaded from
  Cassandra after a restart, but the 4-hour windows and alert episodes start again. After
  the outage test, conditions that were already open were raised again: P014 has 7 alert
  episodes under v1 and 5 under v2 over the same days, with identical scores. Spark's own
  state store would keep them.
  \item \textbf{Replay sees fewer labs.} Because the labs topic is compacted, see
  Section~\ref{sec:decision}.
  \item \textbf{Simulated data and untuned alerts.} The physiology is realistic enough
  to test the pipeline, but the alert rules have not been checked against real
  patients. On our simulated ward they still fire for some stable patients (about 600
  alert episodes over six simulated days). An episode also closes as soon as a score dips
  below its threshold, so a patient near a threshold gets repeated alerts (P031 under
  v2). A margin before closing an episode, agreed with clinicians, would fix both.
  \item \textbf{No tracing and no security.} There is no login, no encryption and no
  audit trail, all of which patient data would need.
\end{itemize}

% =====================================================================
\section{Conclusion}
% =====================================================================

We built a Kappa pipeline for a hospital ward. One Spark stream scores every reading
with one copy of the NEWS2 rule, joins the daily lab results, and stores the result in
Cassandra for a live ward list, an API and a daily report. Kafka keeps the history, and
changing the scoring rule is a replay of that history through the same code. We checked
every result against an answer we knew in advance, and the system met them: P014 is
caught, P007 is not flagged, and only COPD patients change under the new rule.

\begin{thebibliography}{9}
\bibitem{rcp} Royal College of Physicians. \emph{National Early Warning Score (NEWS) 2:
Standardising the assessment of acute-illness severity in the NHS.} London, 2017.
\bibitem{kreps} J. Kreps. ``Questioning the Lambda Architecture.'' O'Reilly Radar, 2014.
\bibitem{kleppmann} M. Kleppmann. \emph{Designing Data-Intensive Applications.}
O'Reilly, 2017.
\bibitem{spark} Apache Software Foundation. \emph{Structured Streaming Programming
Guide}, Spark 3.5.
\end{thebibliography}

\appendix

% =====================================================================
\section{Individual contributions}\label{app:contrib}
% =====================================================================

All three members took part in design discussions, testing and reviewing the report.
The major parts each person built are:

\begin{table}[H]\centering\small
\begin{tabularx}{\linewidth}{|l|X|}
\hline
\rowcolor{tablehead}\textbf{Member} & \textbf{Major parts} \\ \hline
\authorA & The NEWS2 scoring rule and its tests; the physiology simulator and the three
scripted patients; the Spark stream (checking, windows, scoring, lab join, alert
episodes, dead-letter and late paths); the metrics and Prometheus alert rules; writing
this report. \\ \hline
\authorB & The Docker Compose platform (Kafka, Schema Registry, topics); the Cassandra
data model and data access code; the Airflow workflows; the FastAPI service; the replay
stream and the version comparison; the Grafana dashboards; the architecture diagrams. \\ \hline
\authorC & The architecture decision records; the daily PDF report; the test suites for
the API, the replay and the observability configuration; the screenshots and the
measurements used in this report. \\ \hline
\end{tabularx}
\end{table}

% =====================================================================
\section{How to run it}
% =====================================================================

\begin{verbatim}
make clean && make up     start everything from nothing
make topics               list topics with their cleanup policies
make keys                 count patient keys on the vitals topic (40)
make replay               start the v2 stream from offset 0 and wait
make diff                 compare v1 and v2
make cutover VERSION=v2   switch the version the API serves
make chaos-stop-stream    stop the stream to test the alert
make test                 run the automated tests
\end{verbatim}

Ports: Kafka UI 8180, API 8100, Airflow 8182 (admin, admin), Grafana 3100, Prometheus
9190, Alertmanager 9193, Spark UI 4140 (replay 4141), Cassandra 9142.

% =====================================================================
\section{All screenshots}\label{app:shots}
% =====================================================================

\shot{kafka-topics}{0.95}{Kafka UI: the six topics.}
\shot{kafka-labs-settings}{0.95}{Kafka UI: \texttt{labs.results.v1} is compacted.}
\shot{kafka-vitals-messages}{0.95}{Kafka UI: readings on \texttt{vitals.readings.v1},
decoded from Avro through the Schema Registry, keyed by patient.}
\shot{kafka-dlq-messages}{0.95}{Kafka UI: dead-lettered readings with their rejection
reason.}
\shot{kafka-schemas}{0.95}{Kafka UI: the four Avro schemas in the Schema Registry.}
\shot{spark-streaming-v1}{0.95}{Spark UI: the live stream's input rate and batch
duration.}
\shot{spark-streaming-v2}{0.95}{Spark UI: the replay stream catching up.}
\shot{spark-jobs}{0.95}{Spark UI: completed jobs.}
\shot{cassandra-ward-snapshot}{0.95}{Cassandra: the ward snapshot partition, highest
risk first.}
\shot{cassandra-p031-versions}{0.95}{Cassandra: P031 scored under v1 and v2.}
\shot{api-docs}{0.95}{FastAPI documentation page.}
\shot{airflow-dags}{0.95}{Airflow: the five workflows.}
\shot{airflow-lab-ingest-grid}{0.95}{Airflow: \texttt{ward\_lab\_ingest}, one run per
simulated day. The failed runs are the scripted bad days: the day-4 file was quarantined
because of a corrupted reference range, and the day-5 file never arrived.}
\shot{airflow-daily-report-grid}{0.95}{Airflow: \texttt{ward\_daily\_report} runs.}
\shot{airflow-replay-graph}{0.95}{Airflow: \texttt{ward\_replay} stops at the approval
step.}
\shot{prometheus-targets}{0.95}{Prometheus: every scrape target up.}
\shot{alertmanager}{0.95}{Alertmanager: \texttt{WardMonitoringSilent} routed to the
ward and to the on-call engineer.}
\shot{grafana-pipeline-health}{1.0}{Grafana: pipeline health.}
\shot{grafana-pipeline-outage}{1.0}{Grafana: pipeline health during the outage test.}

\end{document}
